\section{Discussion}\label{sec:discussion}
The proof begins with the highest rational band, where convexity and Chambers' formula provide a contour with a large discriminant. The stable Floquet line then gives localization per physical site. That localization is what makes translations, Gram summation and hard frame cutoffs controllable as the denominator grows.

The orbital term moves the common energy origin of the wells. It does not change the count between consecutive clusters. Its Wannier form is a Brillouin-zone average; it is not, in general, its value at Bloch momentum zero. The off-diagonal first-order terms are suppressed by tunnelling. The remaining comparison has two independent relative errors. The two caps on $D$ control those errors separately. Keeping them separate is essential near $b=L/2$.

The final integer count is the bridge from a local well calculation to an irrational gap. Correct covariance gives one symbol in every well. The common Mathieu windows give the same number $D$ in both comparison operators. Their equality forces an exact count $KD$ in each box. No inference from two agreeing numerical box spectra is used in this passage.

The theorem gives a lower bound along a sequence of labels. The endpoint $\beta=L/2$, bounds for all large labels and a matching dynamical law remain open. The change of exponent at $3L/2$ records the precision of the quadratic comparison.

Below the threshold the relative tunnelling scale replaces the small-curvature Mathieu spacing. The continuation is periodic at both far ends and retains only the selected physical well. Its exact covariance makes the full diagonal constant common to all offsets. Keeping the self-energy and estimating its first phase derivative is sufficient for the quadratic model hypothesis; no cubic Taylor assertion across frame jumps is needed.

\subsection*{How this work was done}
The authors developed the results and proofs with assistance from AI models from more than one developer. They take responsibility for every claim. The authors chose the question, directed the work and decided which lines to pursue; the models did most of the exploration, derivation and drafting. Both authors contributed equally.

The source work proceeded through separate research sessions on gap sizes, the arithmetic picture, the polynomial remainder, top-band analyticity and the global reduction. Each was followed by a separate session instructed to refute the submitted argument and inspect its primary sources; these checks are D76--D80 in the working archive, and their repairs take precedence. The initial global-frame argument failed its check; the translated-well construction presented here replaced it. A further refutation-first audit of the assembled manuscript (D81) required repairs to Appendix~\ref{app:mathieu}, to the companion table and to several attributions; this version makes them.

The new below-threshold gap argument, the complementary Green-function estimate and the corrected covariant continuation were developed in separate research sessions, each followed by an independent refutation-first check whose replacement estimates are incorporated here.

Priority and literature searches were AI-assisted; an unsuccessful bounded search is not a priority certificate. No computation is used in any proof.

\label{end:discussion}
